% Generated draft: requires mathematical and source review.
\documentclass[11pt]{article}
\usepackage[margin=25mm]{geometry}
\usepackage{fontspec}
\setmainfont{Latin Modern Roman}
\usepackage{amsmath,amssymb,mathtools,unicode-math}
\setmathfont{Latin Modern Math}
\usepackage{enumitem,needspace,xurl,hyperref}
\setlength{\parindent}{0pt}
\setlength{\parskip}{5pt}
\setlength{\emergencystretch}{3em}
\newcommand{\N}{\mathbb{N}}
\newcommand{\Z}{\mathbb{Z}}
\newcommand{\Q}{\mathbb{Q}}
\newcommand{\R}{\mathbb{R}}
\newcommand{\C}{\mathbb{C}}
\newcommand{\F}{\mathbb{F}}
\newcommand{\A}{\mathbb{A}}
\newcommand{\PP}{\mathbb{P}}
\newcommand{\E}{\mathbb{E}}
\newcommand{\eps}{\varepsilon}
\newcommand{\dd}{\,\mathrm{d}}
\begin{document}
\section*{UnsolvedMath ID: 700218}
\textbf{Source problem code:} catalogue entry
\subsection{Definitions and mathematical statement}
Fix an integer \(n\geq 2\). The supplied problem asks which ring epimorphisms \(f:R\to Q\) induce an epimorphism \(SL_n(R)\to SL_n(Q)\), and, in particular, which rings \(R\) satisfy \(SL_n(R)=E_n(R)\). The source does not specify the class of rings, the convention for ring homomorphisms, or the meaning of ring epimorphism. These qualifications remain unresolved.

For a precise formulation under the additional convention that rings are commutative with identity and homomorphisms preserve identity, define
\[
SL_n(R)=\{A\in M_n(R):\det A=1_R\}.
\]
For \(i\ne j\), let \(e_{ij}\) be the matrix unit, and define the elementary subgroup by
\[
E_n(R)=\left\langle I_n+r e_{ij}:r\in R,\ 1\leq i,j\leq n,\ i\ne j\right\rangle\leq SL_n(R).
\]
A homomorphism \(f:R\to Q\) induces
\[
f_n:SL_n(R)\longrightarrow SL_n(Q),\qquad
(a_{ij})\longmapsto(f(a_{ij})).
\]
Under this convention, the first question is to characterize the ring epimorphisms \(f\) for which
\[
\forall B\in SL_n(Q)\ \exists A\in SL_n(R)\quad f_n(A)=B.
\]
Here an epimorphism of groups means a surjective group homomorphism. The admissible ring maps still depend on an unresolved source convention: a surjective ring homomorphism need not mean the same thing as a categorical ring epimorphism. In the category of commutative rings with identity and identity-preserving homomorphisms, the latter means that, for every ring \(T\) and every pair of homomorphisms \(g,h:Q\to T\),
\[
g\circ f=h\circ f\quad\Longrightarrow\quad g=h.
\]
The source does not determine which of these meanings is intended.

The particular question, under the same commutative-ring convention, is to characterize the rings \(R\) such that every \(A\in SL_n(R)\) is a finite product
\[
A=\prod_{k=1}^{m}(I_n+r_k e_{i_kj_k}),
\qquad m\geq0,\quad r_k\in R,\quad 1\leq i_k,j_k\leq n,\quad i_k\ne j_k,
\]
where the empty product is \(I_n\). The integer \(n\) remains fixed throughout; no simultaneous assertion for all \(n\) is imposed. These definitions provide a conditional commutative-ring formulation, not a determination of the source's intended ring domain. If noncommutative rings are intended, the meaning of \(SL_n\) must also be specified.
\subsection{Short English statement}
For a fixed \(n\geq2\), characterize the ring epimorphisms for which every determinant-one matrix over the target lifts to a determinant-one matrix over the source. Also characterize the rings whose determinant-one matrices are all products of elementary matrices. The determinant formulation assumes commutative rings with identity; the source leaves the ring domain and the meaning of ring epimorphism unspecified.
\subsection{Sources}
\begin{itemize}
\item Supplied local UnsolvedMath registry record: permanent ID \(700218\), exact source code \texttt{catalogue entry}, title \emph{Kourovka Notebook Problem 4.7}; upstream revision \texttt{d371ee917ac9647fbb06cd6fc529deec6f3dda82}.
\item The supplied attribution is to V. A. Artamonov, proposed in 1973. The statement is introduced as a ``Well-known problem.''
\item Supplied reference: \emph{The Kourovka Notebook (Unsolved Problems in Group Theory)}, 4th issue (1973), Problem 4.7; 21st issue edition (2026), \texttt{arXiv:1401.0300v48}. The collection description names E. I. Khukhro and V. D. Mazurov as editors and describes coverage of Issues 1--21 (1965--2026).
\item The registry labels the problem \texttt{open}. This is a supplied status, not an independently verified mathematical or literature assessment. The supplied reference has not been independently consulted for this draft.
\end{itemize}
\end{document}
